% ==============================================================================
% CAPITOLO 4: VARIABILI ALEATORIE DISCRETE, CONTINUE E VETTORI ALEATORI
% ==============================================================================
\chapter{Variabili Aleatorie Discrete, Continue e Vettori Aleatori}
\label{cap:variabili_aleatorie}

\begin{abstract}
Il presente capitolo opera la transizione concettuale fondamentale dalla teoria degli insiemi astratti alla modellazione quantitativa mediante funzioni a valori reali: le variabili aleatorie. Viene formalizzata la nozione di misurabilità e analizzata la Funzione di Ripartizione (CDF) come descrittore universale sia per leggi discrete che per leggi continue. Si approfondiscono le specificità delle variabili discrete (PMF) e delle variabili assolutamente continue (PDF), chiarendo la natura infinitesima della densità e il significato epistemologico della probabilità puntuale nulla. Nella seconda parte, la teoria viene estesa ai vettori aleatori bivariati: densità congiunte, distribuzioni marginali, leggi condizionate, indipendenza stocastica e la formula analitica del cambio di variabile con fattore Jacobiano di distorsione geometrica.
\end{abstract}

% ==============================================================================
% SEZIONE 4.1: IL CONCETTO DI VARIABILE ALEATORIA E LA CDF
% ==============================================================================
\section{Il Concetto di Variabile Aleatoria e la Funzione di Ripartizione (CDF)}
\label{sec:concetto_va_cdf}

Fino ad ora abbiamo trattato gli esiti sperimentali come elementi qualitativi o insiemistici di uno spazio campionario astratto $\Omega$ (ad esempio, sequenze di teste e croci, colori di biglie o campioni di guasto). Tuttavia, per applicare gli strumenti del calcolo differenziale e integrale, dell'algebra lineare e dell'analisi numerica, è indispensabile tradurre ciascun esito elementare $\omega \in \Omega$ in una grandezza numerica quantitativa reale $x \in \mathbb{R}$.

\begin{definizione}{Variabile Aleatoria Reale}{def_variabile_aleatoria}
Sia $(\Omega, \mathcal{F}, \mathbb{P})$ uno spazio di probabilità. Una \textbf{variabile aleatoria reale} (v.a.) $X$ è una funzione misurabile che mappa lo spazio campionario nella retta reale:
\begin{equation}
    X: \Omega \longrightarrow \mathbb{R}, \qquad \omega \longmapsto X(\omega)
\end{equation}
La condizione di \textbf{misurabilità} impone che per ogni intervallo di Borel $I \subseteq \mathbb{R}$, l'insieme degli esiti elementari la cui immagine ricade in $I$ sia un evento appartenente alla $\sigma$-algebra $\mathcal{F}$:
\begin{equation}
    \{\omega \in \Omega : X(\omega) \in I\} \equiv X^{-1}(I) \in \mathcal{F}
\end{equation}
L'insieme dei valori effettivamente assunti dalla funzione è detto \textbf{supporto} o immagine della variabile aleatoria: $\operatorname{Im}(X) = X(\Omega) \subseteq \mathbb{R}$.
\end{definizione}

Per caratterizzare probabilisticamente il comportamento di $X$ senza dover fare continuamente riferimento allo spazio originario $\Omega$, si introduce la **Funzione di Ripartizione Cumulativa**.

\begin{definizione}{Funzione di Ripartizione Cumulativa (CDF)}{def_cdf}
La \textbf{funzione di ripartizione} (o Cumulative Distribution Function, CDF) di una variabile aleatoria $X$ è la funzione a valori reali $F_X: \mathbb{R} \to [0, 1]$ definita da:
\begin{equation}
    F_X(x) \coloneqq \mathbb{P}(X \le x) = \mathbb{P}(\{\omega \in \Omega : X(\omega) \le x\})
\end{equation}
\end{definizione}

\begin{teorema}{Proprietà Caratteristiche Fondamentali della CDF}{prop_cdf}
Qualsiasi funzione di ripartizione $F_X(x)$ soddisfa inderogabilmente le seguenti quattro proprietà matematiche:
\begin{enumerate}
    \item \textbf{Monotonia non decrescente:} Se $x_1 < x_2$, allora $F_X(x_1) \le F_X(x_2)$ (perché $\{X \le x_1\} \subseteq \{X \le x_2\}$).
    \item \textbf{Limiti agli estremi:}
    \begin{equation}
        \lim_{x \to -\infty} F_X(x) = \mathbb{P}(\emptyset) = 0, \qquad \lim_{x \to +\infty} F_X(x) = \mathbb{P}(\Omega) = 1
    \end{equation}
    \item \textbf{Continuità da destra:} Per ogni $x_0 \in \mathbb{R}$, $F_X(x_0^+) \coloneqq \lim_{h \to 0^+} F_X(x_0 + h) = F_X(x_0)$.
    \item \textbf{Probabilità degli intervalli e ampiezza dei salti:}
    \begin{align}
        \mathbb{P}(a < X \le b) &= F_X(b) - F_X(a) \\
        \mathbb{P}(X < x_0) &= F_X(x_0^-) \coloneqq \lim_{h \to 0^+} F_X(x_0 - h) \\
        \mathbb{P}(X = x_0) &= F_X(x_0) - F_X(x_0^-) = \text{ampiezza del salto di discontinuità in } x_0
    \end{align}
\end{enumerate}
\end{teorema}

\begin{figure}[htbp]
    \centering
    \begin{tikzpicture}[scale=0.85, >=Stealth]
    % -------------------------------------------------------------
    % PANNELLO A: CDF DISCRETA
    % -------------------------------------------------------------
    \begin{scope}[shift={(0,0)}]
        \node[above, font=\bfseries\footnotesize, text=navyblue] at (2.5,4.2) {(a) CDF Discreta $F_X(x)$ (a scalini)};
        \draw[step=1cm, gray!15, very thin] (-0.5,-0.2) grid (5.2,4.0);

        \draw[->, thick, navyblue] (-0.5,0) -- (5.3,0) node[right, font=\footnotesize] {$x$};
        \draw[->, thick, navyblue] (0,-0.3) -- (0,4.2) node[above, font=\footnotesize] {$F_X(x)$};

        % Asintoto y = 1
        \draw[dashed, crimsonred!70] (-0.3,3.5) -- (5.0,3.5) node[right, font=\tiny] {$1$};
        \node[left, font=\tiny] at (0,3.5) {$1$};
        \node[left, font=\tiny] at (0,0) {$0$};

        % Scalini
        % x < 1: F(x) = 0
        \draw[ultra thick, coolblue] (-0.5,0) -- (1,0);
        \draw[coolblue, fill=white, thick] (1,0) circle (2pt);

        % 1 <= x < 2: F(x) = 0.3
        \fill[coolblue] (1,1.05) circle (2pt);
        \draw[ultra thick, coolblue] (1,1.05) -- (2,1.05);
        \draw[coolblue, fill=white, thick] (2,1.05) circle (2pt);
        \draw[dotted, gray] (1,0) -- (1,1.05);

        % 2 <= x < 3.5: F(x) = 0.7
        \fill[coolblue] (2,2.45) circle (2pt);
        \draw[ultra thick, coolblue] (2,2.45) -- (3.5,2.45);
        \draw[coolblue, fill=white, thick] (3.5,2.45) circle (2pt);
        \draw[dotted, gray] (2,0) -- (2,2.45);

        % x >= 3.5: F(x) = 1.0
        \fill[coolblue] (3.5,3.5) circle (2pt);
        \draw[ultra thick, coolblue] (3.5,3.5) -- (5.0,3.5);
        \draw[dotted, gray] (3.5,0) -- (3.5,3.5);

        \node[below, font=\tiny] at (1,0) {$x_1$};
        \node[below, font=\tiny] at (2,0) {$x_2$};
        \node[below, font=\tiny] at (3.5,0) {$x_3$};

        \node[font=\tiny\itshape, text=forestgreen!80!black, anchor=west] at (1.1,1.75) {Salto $= p(x_2)$};
        \draw[->, forestgreen, thick] (1.05,1.75) -- (1.95,1.75);
    \end{scope}

    % -------------------------------------------------------------
    % PANNELLO B: CDF CONTINUA
    % -------------------------------------------------------------
    \begin{scope}[shift={(6.8,0)}]
        \node[above, font=\bfseries\footnotesize, text=navyblue] at (2.5,4.2) {(b) CDF Continua $F_X(x)$ (liscia)};
        \draw[step=1cm, gray!15, very thin] (-0.5,-0.2) grid (5.2,4.0);

        \draw[->, thick, navyblue] (-0.5,0) -- (5.3,0) node[right, font=\footnotesize] {$x$};
        \draw[->, thick, navyblue] (0,-0.3) -- (0,4.2) node[above, font=\footnotesize] {$F_X(x)$};

        \draw[dashed, crimsonred!70] (-0.3,3.5) -- (5.0,3.5) node[right, font=\tiny] {$1$};
        \node[left, font=\tiny] at (0,3.5) {$1$};
        \node[left, font=\tiny] at (0,0) {$0$};

        \draw[ultra thick, forestgreen, smooth] plot coordinates {
            (-0.5, 0.02) (0.5, 0.1) (1.5, 0.5) (2.5, 1.75) (3.5, 3.0) (4.5, 3.45) (5.0, 3.5)
        };

        \draw[dashed, navyblue] (2.5,0) -- (2.5,1.75) -- (0,1.75);
        \node[below, font=\tiny] at (2.5,0) {Mediana};
        \node[left, font=\tiny] at (0,1.75) {$0.5$};
    \end{scope}
\end{tikzpicture}

    \caption{Confronto tipologico tra la CDF di una variabile discreta (a sinistra), caratterizzata da una struttura a scalini con salti finiti nei punti del supporto, e la CDF di una variabile assolutamente continua (a destra), priva di salti e ovunque continua e derivabile.}
    \label{fig:cdf_confronto}
\end{figure}

% ==============================================================================
% SEZIONE 4.2: VARIABILI ALEATORIE DISCRETE E FUNZIONE DI MASSA (PMF)
% ==============================================================================
\section{Variabili Aleatorie Discrete e Funzione di Massa (PMF)}
\label{sec:va_discrete_pmf}

\begin{definizione}{Variabile Aleatoria Discreta e PMF}{def_pmf}
Una variabile aleatoria $X$ si dice \textbf{discreta} se il suo supporto $\operatorname{Im}(X) = \{x_1, x_2, \dots\}$ è un insieme finito o al più numerabile (isomorfo a $\mathbb{N}$). La sua distribuzione è descritta dalla \textbf{funzione di massa di probabilità} (Probability Mass Function, PMF) $p_X(x)$:
\begin{equation}
    p_X(x) \coloneqq \mathbb{P}(X = x)
\end{equation}
La PMF soddisfa i vincoli di positività e normalizzazione:
\begin{equation}
    p_X(x) \ge 0 \quad \forall x \in \operatorname{Im}(X), \qquad \sum_{x \in \operatorname{Im}(X)} p_X(x) = 1
\end{equation}
La CDF associata è una funzione a scalini monotona crescente:
\begin{equation}
    F_X(x) = \sum_{x_k \le x} p_X(x_k)
\end{equation}
\end{definizione}

% ==============================================================================
% SEZIONE 4.3: VARIABILI ALEATORIE CONTINUE E FUNZIONE DI DENSITÀ (PDF)
% ==============================================================================
\section{Variabili Aleatorie Continue e Funzione di Densità (PDF)}
\label{sec:va_continue_pdf}

Nelle grandezze fisiche come tempo, massa, temperatura o pressione, il supporto della variabile aleatoria è un continuo (un intervallo o l'intera retta reale $\mathbb{R}$).

\begin{definizione}{Variabile Aleatoria Assolutamente Continua e PDF}{def_pdf}
Una variabile aleatoria $X$ si dice \textbf{assolutamente continua} se esiste una funzione integrabile e non negativa $f_X: \mathbb{R} \to [0, \infty)$, detta \textbf{funzione di densità di probabilità} (Probability Density Function, PDF), tale che:
\begin{equation}
    F_X(x) = \mathbb{P}(X \le x) = \int_{-\infty}^x f_X(t) \, dt \qquad (\forall x \in \mathbb{R})
\end{equation}
Una PDF ammissibile deve soddisfare le due condizioni:
\begin{equation}
    f_X(x) \ge 0 \quad \forall x \in \mathbb{R}, \qquad \int_{-\infty}^{+\infty} f_X(x) \, dx = 1
\end{equation}
Nei punti in cui $f_X$ è continua, per il Teorema Fondamentale del Calcolo Integrale si ha:
\begin{equation}
    f_X(x) = \frac{d}{dx} F_X(x) = F_X'(x)
\end{equation}
\end{definizione}

\begin{figure}[htbp]
    \centering
    \begin{tikzpicture}[scale=0.85, >=Stealth]
    \draw[->, thick, navyblue] (-0.5,0) -- (7.5,0) node[right, font=\footnotesize] {$x$};
    \draw[->, thick, navyblue] (0,-0.3) -- (0,4.0) node[above, font=\footnotesize] {$f_X(x)$ (PDF)};

    % Curva PDF continua (Gaussiana-like)
    \draw[domain=0.5:6.5, smooth, variable=\x, very thick, navyblue]
        plot ({\x}, {3.5 * exp(-(\x-3.5)*(\x-3.5)/2.2)});

    % Area sottesa tra a = 2.2 e b = 4.8
    \fill[coolblue!30, domain=2.2:4.8, variable=\x]
        (2.2, 0) --
        plot ({\x}, {3.5 * exp(-(\x-3.5)*(\x-3.5)/2.2)}) --
        (4.8, 0) -- cycle;

    \draw[thick, dashed, navyblue] (2.2,0) -- (2.2, {3.5 * exp(-(2.2-3.5)*(2.2-3.5)/2.2)});
    \draw[thick, dashed, navyblue] (4.8,0) -- (4.8, {3.5 * exp(-(4.8-3.5)*(4.8-3.5)/2.2)});

    \node[below, font=\small\bfseries] at (2.2,0) {$a$};
    \node[below, font=\small\bfseries] at (4.8,0) {$b$};

    \node[font=\small\bfseries, text=navyblue, align=center] at (3.5, 1.2) {
        $\mathbb{P}(a \le X \le b) = \displaystyle\int_a^b f_X(x) \, dx$
    };

    % Striscia infinitesima
    \draw[fill=warmamber!60, draw=warmamber!90!black] (5.4,0) rectangle (5.55, {3.5 * exp(-(5.45-3.5)*(5.45-3.5)/2.2)});
    \node[above right, font=\tiny, text=warmamber!90!black] at (5.5, 0.8) {$f_X(x)dx \approx \mathbb{P}(x \le X \le x+dx)$};
\end{tikzpicture}

    \caption{Significato geometrico della funzione di densità $f_X(x)$. La probabilità che $X$ assuma valori nell'intervallo infinitesimo $[x, x+dx]$ è l'area del rettangolo elementare $f_X(x)dx$. La probabilità dell'intervallo $[a, b]$ corrisponde all'area totale sottesa dalla curva tra $a$ e $b$.}
    \label{fig:pdf_area}
\end{figure}

\paragraph{Il Paradosso della Probabilità Puntuale Nulla nel Continuo: $\mathbb{P}(X = x_0) = 0$.}
Un aspetto che genera frequente smarrimento è che, per una variabile continua, la probabilità che $X$ assuma un qualsiasi valore puntuale prefissato $x_0 \in \mathbb{R}$ è **rigorosamente pari a zero**:
\begin{equation}
    \mathbb{P}(X = x_0) = \int_{x_0}^{x_0} f_X(t) \, dt = 0
\end{equation}
\emph{Come può un evento avere probabilità zero senza essere impossibile?}
L'analogia fisica chiarificatrice è quella della densità di massa di un corpo continuo. Un'asta metallica di massa totale pari a $1\text{ kg}$ possiede una densità lineare $\rho(x) \text{ kg/m}$. Un singolo punto geometrico $x_0$ ha volume nullo e dunque possiede massa rigorosamente zero ($dm = 0$). Tuttavia, l'asta è formata dall'unione di infiniti punti. 

Nel continuo, il valore $f_X(x_0)$ non è una probabilità, ma una **densità di probabilità locale per unità di lunghezza** (misurata in $[\text{unità}]^{-1}$). La probabilità fisica non risiede mai nei singoli punti isolati, bensì nelle aree integrate su intervalli:
\begin{equation}
    \mathbb{P}(a \le X \le b) = \mathbb{P}(a < X \le b) = \mathbb{P}(a \le X < b) = \mathbb{P}(a < X < b) = \int_a^b f_X(x) \, dx
\end{equation}

% ==============================================================================
% SEZIONE 4.4: VETTORI ALEATORI BIVARIATI
% ==============================================================================
\section{Vettori Aleatori Bivariati, Distribuzioni Congiunte e Marginali}
\label{sec:vettori_aleatori}

\begin{definizione}{Vettore Aleatorio Continuo e PDF Congiunta}{def_vettore_continuo}
Una coppia di variabili aleatorie $(X, Y)$ definite sul medesimo spazio di probabilità forma un \textbf{vettore aleatorio bivariato in $\mathbb{R}^2$}. Se ammette una funzione di densità congiunta $f_{X,Y}(x, y) \ge 0$, la probabilità che il punto casuale $(X, Y)$ cada in una regione bidimensionale $A \subseteq \mathbb{R}^2$ è data dall'integrale doppio:
\begin{equation}
    \mathbb{P}((X, Y) \in A) = \iint_A f_{X,Y}(x, y) \, dx \, dy
\end{equation}
con la condizione di normalizzazione sul piano $\int_{-\infty}^{+\infty}\int_{-\infty}^{+\infty} f_{X,Y}(x, y) \, dx \, dy = 1$.
\end{definizione}

\subsection{Densità Marginali e Densità Condizionate}

\begin{definizione}{Densità Marginali}{def_marginali}
Integrando la densità congiunta rispetto ad una delle due variabili su tutto il suo supporto, si ottiene la \textbf{densità marginale} della variabile residua (principio di saturazione dello spazio):
\begin{equation}
    f_X(x) = \int_{-\infty}^{+\infty} f_{X,Y}(x, y) \, dy, \qquad f_Y(y) = \int_{-\infty}^{+\infty} f_{X,Y}(x, y) \, dx
\end{equation}
\end{definizione}

\begin{definizione}{Densità Condizionata}{def_densita_condizionata}
Fissato un valore $x$ tale che $f_X(x) > 0$, la \textbf{funzione di densità condizionata} di $Y$ dato $X = x$ è definita da:
\begin{equation}
    f_{Y \mid X}(y \mid x) = \frac{f_{X,Y}(x, y)}{f_X(x)}
\end{equation}
Essa rappresenta il profilo di sezione geometrica della superficie $f_{X,Y}(x, y)$ lungo la retta $X = x$, rinormalizzato affinché l'area sottesa lungo $y$ sia esattamente unitaria: $\int_{-\infty}^{+\infty} f_{Y \mid X}(y \mid x) \, dy = 1$.
\end{definizione}

\subsection{Indipendenza tra Variabili Aleatorie}

\begin{definizione}{Indipendenza di Variabili Aleatorie}{def_indipendenza_va}
Due variabili aleatorie $X$ e $Y$ si dicono \textbf{stocasticamente indipendenti} ($X \independent Y$) se e solo se la funzione di ripartizione congiunta fattorizza nel prodotto delle CDF marginali per ogni coppia $(x, y) \in \mathbb{R}^2$:
\begin{equation}
    F_{X,Y}(x, y) = F_X(x) \cdot F_Y(y)
\end{equation}
Nel caso assolutamente continuo, ciò equivale alla fattorizzazione quasi ovunque della densità congiunta:
\begin{equation}
    f_{X,Y}(x, y) = f_X(x) \cdot f_Y(y) \iff f_{Y \mid X}(y \mid x) = f_Y(y)
\end{equation}
\end{definizione}

% ==============================================================================
% SEZIONE 4.5: TRASFORMAZIONI DI VARIABILI E FATTORE JACOBIANO
% ==============================================================================
\section{Trasformazioni di Variabili Aleatorie e Metodo dello Jacobiano}
\label{sec:trasformazioni_va}

In molte applicazioni ingegneristiche si conosce la densità di una variabile di ingresso $X$ e si desidera calcolare la densità della grandezza di uscita $Y = g(X)$ (ad esempio, tensione ai capi di un diodo non lineare $V = R I^2$).

\begin{teorema}{Trasformazione Monotona 1D}{thm_trasformazione_1d}
Sia $X$ una v.a. continua con densità $f_X(x)$ e supporto $(a, b)$. Sia $Y = g(X)$ una trasformazione differenziabile e \textbf{strettamente monotona} (quindi invertibile con inversa $x = g^{-1}(y)$). Allora la densità di $Y$ è data da:
\begin{equation}
    f_Y(y) = f_X(g^{-1}(y)) \cdot \left| \frac{d}{dy} g^{-1}(y) \right| = \frac{f_X(x)}{|g'(x)|} \Bigg|_{x = g^{-1}(y)}
\end{equation}
\end{teorema}

\begin{dimostrazione}[Dimostrazione per il caso strettamente crescente $g'(x) > 0$]
Se $g$ è strettamente crescente, l'evento $\{Y \le y\}$ equivale a $\{g(X) \le y\} = \{X \le g^{-1}(y)\}$.
Pertanto la CDF di $Y$ è:
\begin{equation}
    F_Y(y) = \mathbb{P}(Y \le y) = \mathbb{P}(X \le g^{-1}(y)) = F_X(g^{-1}(y))
\end{equation}
Derivando rispetto a $y$ mediante la regola della catena (Teorema di derivazione delle funzioni composte):
\begin{equation}
    f_Y(y) = \frac{d}{dy} F_Y(y) = F_X'(g^{-1}(y)) \cdot \frac{d}{dy} g^{-1}(y) = f_X(g^{-1}(y)) \cdot \frac{d}{dy} g^{-1}(y)
\end{equation}
Nel caso decrescente compare un segno meno che viene assorbito dal valore assoluto $|\frac{d}{dy}g^{-1}(y)|$. $\blacksquare$
\end{dimostrazione}
